\chapter{Resultados experimentales}
\label{chap:resultados}

Este capítulo presenta, primero, la caracterización del sistema de control y adquisición
(temporización y repetibilidad) y, después, los resultados sobre el mecanismo: la evolución de la
tensión del cable y del par del motor a lo largo del recorrido, la separación de sus componentes y
la comparación con el modelo. Todas las figuras y cifras se obtienen con
\fichero{memoria/figuras/generar\_figuras.py} a partir de los ficheros de \fichero{cycles/}, sin
ninguna edición manual.

\section{Condiciones de ensayo}

\completar{fotografía del prototipo montado, con la caracola, el cable, la célula de carga, la barra,
la masa y los finales de carrera identificados}

Los ensayos se realizaron los días 14 y 15 de septiembre de 2026 con la masa de \SI{2}{\kilo\gram}.
Antes de cada serie se inicializó el accionamiento y se ejecutó el \emph{homing}. Cada ensayo es un
ciclo automático que parte del origen (A), avanza a velocidad constante hasta el final de carrera B
y vuelve al origen. Los parámetros del accionamiento fueron los de la tabla~\ref{tab:parametros}
(rigidez 10, rampas 5) \verificar{confirmar que no se modificaron durante la campaña}.

Se grabaron 15 ficheros (tabla~\ref{tab:ensayos}). Uno de ellos, a \SI{300}{rpm}, se interrumpió
a los \SI{1.15}{\second} y se descarta. Quedan 14 ensayos válidos con 2744 muestras: 12 a
\SI{30}{rpm}, que suman 14 ciclos porque dos de ellos tienen dos repeticiones, uno a \SI{20}{rpm} y
otro a \SI{40}{rpm}.

\begin{longtable}{@{}lcccccl@{}}
\caption{Ensayos grabados. Velocidad de consigna, repeticiones, muestras, duración y posición
máxima alcanzada (final de carrera B).}\label{tab:ensayos}\\
\toprule
Fichero & rpm & Rep. & Muestras & $t$ (s) & $p_{\max}$ (rev) & Observaciones \\
\midrule
\endfirsthead
\toprule
Fichero & rpm & Rep. & Muestras & $t$ (s) & $p_{\max}$ (rev) & Observaciones \\
\midrule
\endhead
20260914-183121 & 20 & 1 & 252 & 12,55 & 1,763 & \\
20260914-184729 & 30 & 2 & 332 & 16,55 & 1,763 & \\
20260914-190535 & 30 & 2 & 340 & 16,95 & 1,762 & \\
20260914-190826 & 31 & 1 & 160 & 7,95 & 1,755 & \\
20260914-194105 & 40 & 1 & 128 & 6,35 & 1,761 & \\
20260914-194640 & 300 & 1 & 24 & 1,15 & 1,810 & descartado \\
20260914-194648 & 30 & 1 & 176 & 8,75 & 1,717 & inicio en $-0{,}10$ \\
20260914-195242 & 30 & 1 & 168 & 8,35 & 1,755 & \\
20260914-195540 & 30 & 1 & 172 & 8,55 & 1,760 & \\
20260914-202156 & 30 & 1 & 168 & 8,35 & 1,752 & \\
20260914-202425 & 30 & 1 & 168 & 8,36 & 1,758 & \\
20260914-202801 & 30 & 1 & 168 & 8,35 & 1,750 & \\
20260914-203733 & 30 & 1 & 172 & 8,55 & 1,756 & \\
20260915-174420 & 30 & 1 & 168 & 8,35 & 1,782 & \\
20260915-181346 & 30 & 1 & 172 & 8,55 & 1,759 & ensayo tipo \\
\bottomrule
\end{longtable}

Los ensayos se grabaron con una versión del servidor que todavía no escribía la calibración en la
cabecera (\codigo{lc1\_cuentas\_por\_N=0}). Como los ficheros contienen las cuentas brutas, los
valores de fuerza de este capítulo se obtienen aplicando después la calibración provisional de la
interfaz: \num{14.41} cuentas/N y \SI{51.78}{\newton} en el punto de tara, con
$F = \text{lc1\_base}/14{,}41 + 51{,}78$ (sección~\ref{sec:res-limitaciones}).

\section{Temporización del lazo de control}
\label{sec:res-temporizacion}

La figura~\ref{fig:temporizacion} muestra la distribución del periodo entre muestras consecutivas y
del tiempo de trabajo de cada vuelta del lazo, medidos por el propio firmware. En las 2744 muestras:

\begin{itemize}
  \item el periodo medio es \SI{50.004}{\milli\second}, con una desviación típica de
        \SI{0.066}{\milli\second}. El \SI{99.56}{\percent} de los periodos son exactamente de
        \SI{50}{\milli\second} y el resto de \SI{51}{\milli\second}, el grano de \codigo{millis()}.
        La frecuencia de muestreo es de \SI{20.0}{\hertz} y cumple RNF1 con holgura;
  \item el tiempo de trabajo del lazo está entre \num{10.82} y \SI{11.02}{\milli\second}, con media
        de \SI{10.97}{\milli\second}. Es el \SI{22}{\percent} del periodo y deja
        \SI{39}{\milli\second} de margen. Su dispersión, de sólo \SI{0.2}{\milli\second}, confirma
        que el lazo está dominado por las tres transacciones Modbus, de duración muy regular
        (sección~\ref{sec:fw-lazo}).
\end{itemize}

\begin{figure}[tbp]
\centering
\includegraphics[width=\textwidth]{temporizacion.pdf}
\caption{Izquierda: periodo entre muestras (escala logarítmica). Derecha: tiempo de trabajo por
vuelta del lazo. Todos los ensayos válidos.}
\label{fig:temporizacion}
\end{figure}

La latencia de parada (RNF2) está acotada por la estructura del lazo. Un final de carrera se lee y
la consigna nula se escribe en la misma vuelta, así que la reacción es como máximo de un periodo
más la rampa de deceleración del accionamiento. En el ensayo tipo, el final B se detecta en
$p=\SI{1.737}{rev}$ y el eje se detiene en \SI{1.759}{rev}: un sobrerrecorrido de
\SI{22}{mrev}, coherente con recorrer a \SI{30}{rpm} durante un periodo
(\SI{25}{mrev}).

\section{Ensayo tipo}

La figura~\ref{fig:ensayo-tipo} muestra el ensayo del 15 de septiembre a las 18:13:46. La secuencia
es la siguiente:

\begin{enumerate}
  \item El eje arranca con el final A pulsado (sombreado rojo), que se libera a los
        \SI{0.2}{\second}. La velocidad alcanza la consigna en unos \SI{0.3}{\second}.
  \item Hacia B, que corresponde a bajar la masa, la tensión del cable crece de forma continua. El
        par del motor se mantiene prácticamente constante hasta $p\approx\SI{1.3}{rev}$ y después
        disminuye. En ese último tramo la velocidad real cae a unas \SI{20}{rpm}, pese a que la
        consigna sigue en \SI{30}{rpm}.
  \item A los \SI{4.1}{\second} se activa el final B (sombreado azul). El firmware invierte la
        consigna y el eje vuelve hacia A, subiendo la masa, con un par mayor que a la ida por efecto
        de las pérdidas.
  \item Al llegar a la posición $p\le\SI{0.02}{rev}$ el ciclo termina. La posición final es
        \SI{-9.5}{mrev} y la duración total \SI{8.55}{\second}.
\end{enumerate}

La lectura bruta de la célula aparece en escalones de \SI{200}{\milli\second} porque sólo se
transmite en la cabecera de cada lote. La lectura filtrada, que viaja en cada muestra, va retrasada
respecto a ella por la media móvil de ocho lecturas.

\begin{figure}[tbp]
\centering
\includegraphics[width=\textwidth]{ensayo_tipo.pdf}
\caption{Ensayo tipo a \SI{30}{rpm}: posición, velocidad, célula de carga (con la escala provisional
en newtons a la derecha) y par del motor.}
\label{fig:ensayo-tipo}
\end{figure}

\section{Repetibilidad de posición}

La figura~\ref{fig:repetibilidad} resume los 14 ciclos a \SI{30}{rpm}.

\begin{itemize}
  \item \textbf{Posición del final de carrera B}: media \SI{1.7553}{rev} y desviación típica
        \SI{13.4}{mrev}, con valores entre \num{1.7172} y \SI{1.7815}{rev}. El valor más
        bajo corresponde al ensayo que arrancó a \SI{-0.10}{rev}; sin él, el rango se reduce a
        \SIrange{1.750}{1.782}{rev}. Esta dispersión combina la repetibilidad del propio final de
        carrera con la deriva de la posición integrada a lo largo de 1,75 revoluciones, así que es
        una cota superior de ambas: menos del \SI{1}{\percent} del recorrido.
  \item \textbf{Posición final en A}: media \SI{+0.6}{mrev}, desviación típica
        \SI{10.7}{mrev} y máximo absoluto \SI{16.6}{mrev}. Siempre está dentro
        de la tolerancia de parada de $\pm\SI{20}{mrev}$. Esta cifra mide la precisión de
        la parada respecto a la posición integrada, no respecto al mecanismo, porque es la propia
        posición integrada la que decide la parada.
\end{itemize}

Para dar escala: con la reducción $N=4$, \SI{10}{mrev} de motor equivalen a
\SI{0.9}{\degree} de giro de la caracola.

\begin{figure}[tbp]
\centering
\includegraphics[width=\textwidth]{repetibilidad.pdf}
\caption{Repetibilidad en los 14 ciclos a \SI{30}{rpm}. Izquierda: posición a la que se alcanza B.
Derecha: posición final en A, con la banda de tolerancia de $\pm\SI{20}{mrev}$.}
\label{fig:repetibilidad}
\end{figure}

\section{Tensión del cable y par del motor frente a la posición}

La figura~\ref{fig:fuerza-par-pos} superpone todas las muestras de los ensayos a \SI{30}{rpm} en las
que el eje se mueve de forma estable, es decir, con velocidad del mismo signo que la consigna y al
menos el \SI{60}{\percent} de su valor. Se aprecian tres hechos:

\begin{enumerate}
  \item Los ensayos son muy repetibles: las nubes de puntos de los distintos ciclos se superponen.
  \item La tensión crece de forma monótona desde el origen hasta B, lo que confirma que ir hacia B
        baja la masa (el cable se alarga y, por la ecuación~\eqref{eq:tension}, la tensión aumenta).
  \item El par del motor es mayor hacia A (subida) que hacia B (bajada) en todo el recorrido, como
        corresponde a unas pérdidas que se oponen siempre al movimiento. Su valor es casi constante
        hasta \SI{1.3}{rev} y cae en el tramo final.
\end{enumerate}

\begin{figure}[tbp]
\centering
\includegraphics[width=\textwidth]{fuerza_par_posicion.pdf}
\caption{Tensión del cable (arriba) y par del motor (abajo) frente a la posición, en los 12 ensayos a
\SI{30}{rpm}. Azul: hacia B (baja la masa). Rojo: hacia A (sube la masa).}
\label{fig:fuerza-par-pos}
\end{figure}

\section{Separación de componentes y efecto de la caracola}
\label{sec:res-caracola}

Para aplicar el método de la sección~\ref{sec:mod-separacion}, las muestras se agrupan en intervalos
de posición de \SI{0.05}{rev} y en cada uno se promedian por separado la fuerza y el par de cada
sentido. Sólo se usan los intervalos con al menos ocho muestras en ambos sentidos, que resultan 34
entre \num{0.03} y \SI{1.73}{rev}. La figura~\ref{fig:descomposicion} muestra el resultado.

\begin{figure}[tbp]
\centering
\includegraphics[width=\textwidth]{descomposicion_par.pdf}
\caption{Arriba: tensión media por intervalo en cada sentido y su promedio. Centro: par en cada
sentido, componente gravitatoria $\tau_g$ y de pérdidas $\tau_f$. Abajo: radio efectivo relativo
$\tau_g/F$, normalizado a su máximo.}
\label{fig:descomposicion}
\end{figure}

\paragraph{Tensión.} Promediando ambos sentidos, la tensión pasa de \SI{30.6}{\newton} en el origen a
\SI{42.8}{\newton} en B. Las curvas de ida y vuelta difieren en promedio sólo
\SI{0.27}{\newton}: una vez agrupadas por posición, la histéresis aparente que introduce el filtro de
la célula prácticamente desaparece.

\paragraph{Pérdidas.} La componente de pérdidas es casi constante a lo largo del recorrido, con un
valor medio de \SI{4.9}{\uacc}, el \SI{14}{\percent} de la componente gravitatoria. Una pérdida
independiente de la posición y del sentido corresponde a un rozamiento de tipo Coulomb en la
transmisión y en el guiado del cable.

\paragraph{Compensación en el tramo central.} Entre \num{0.3} y \SI{1.2}{rev}, lejos de las
aceleraciones del arranque y del final de carrera:

\begin{center}
\begin{tabular}{@{}lccc@{}}
\toprule
Magnitud & Mínimo & Máximo & Variación relativa \\
\midrule
Tensión del cable $F$ & \SI{31.9}{\newton} & \SI{37.8}{\newton} & \SI{17.0}{\percent} \\
Par gravitatorio del motor $\tau_g$ & \SI{36.4}{\uacc} & \SI{38.9}{\uacc} & \SI{6.4}{\percent} \\
\bottomrule
\end{tabular}
\end{center}

Con un tambor cilíndrico, el par sería proporcional a la tensión (ecuación~\eqref{eq:par}) y variaría
también un \SI{17}{\percent}. La caracola reduce esa variación a un \SI{6.4}{\percent}: compensa del
orden del \SI{60}{\percent} de la variación de la carga. El resultado sólo usa variaciones relativas,
así que no depende de la escala desconocida del par, y tampoco de la calibración absoluta de la
célula, salvo en el valor de la fuerza en la tara, que fija el denominador. \textbf{Este es el
resultado experimental central del trabajo.}

\paragraph{Tramo final.} A partir de \SI{1.3}{rev} el par gravitatorio cae hasta \SI{18}{\uacc}
mientras la tensión sigue creciendo, y el radio efectivo baja al \SI{34}{\percent} de su valor
inicial. Hay tres explicaciones posibles, que no se han podido discriminar con estos ensayos: (a) el
perfil real de la caracola se estrecha bruscamente al final de la espiral; (b) el cable llega al
extremo de la espiral y cambia su punto de salida; (c) la caída de velocidad observada en ese tramo
(figura~\ref{fig:ensayo-tipo}) rompe la hipótesis de régimen cuasiestático e igual velocidad en ambos
sentidos \verificar{medir el perfil de radio de la caracola real y comprobar la salida del cable al
final del recorrido}.

\section{Influencia de la velocidad}

La figura~\ref{fig:velocidad} compara el par en ensayos a \num{20}, \num{30} y \SI{40}{rpm}. Los
pares medios son \SIlist{38.7;38.4;39.7}{\uacc} hacia A y \SIlist{30.2;27.6;28.9}{\uacc} hacia B,
respectivamente. No se aprecia una tendencia con la velocidad: las diferencias, inferiores al
\SI{10}{\percent}, están en el orden de la variación entre ensayos. Esto confirma que en este rango
las pérdidas no son viscosas y que el ensayo es cuasiestático. La posición de B coincide en los tres
casos con una dispersión de \SI{4}{mrev} (\SIlist{1.763;1.759;1.761}{rev}).

\begin{figure}[tbp]
\centering
\includegraphics[width=\textwidth]{efecto_velocidad.pdf}
\caption{Par del motor frente a la posición a \num{20}, \num{30} y \SI{40}{rpm}. Línea continua:
hacia B. Discontinua: hacia A.}
\label{fig:velocidad}
\end{figure}

\section{Coherencia con el modelo}

Aunque el modelo se ha evaluado con cotas nominales, conviene comprobar que los órdenes de magnitud
son coherentes:

\begin{itemize}
  \item La tensión medida (\SIrange{30.6}{42.8}{\newton}) está dentro del rango teórico con la masa
        de \SI{2}{\kilo\gram} (\SIrange{28.5}{56.2}{\newton}).
  \item Invirtiendo $T=Kc$, esa tensión corresponde a un cable de \SIrange{209.6}{293.2}{\milli\metre}
        y, por la ecuación~\eqref{eq:angulo}, a una barra que pasa de unos \SI{63}{\degree} a unos
        \SI{7}{\degree}.
  \item Esos \SI{83.6}{\milli\metre} de cable se pagan en \SI{1.755}{rev} de motor, es decir,
        \SI{0.439}{rev} de caracola con $N=4$. El radio medio de enrollamiento resulta de
        \SI{30.3}{\milli\metre}, un valor razonable para la pieza.
  \item Con ese radio medio y $F\approx\SI{35}{\newton}$, el par gravitatorio en el motor sería de
        unos \SI{0.27}{\newton\metre} (con $\eta = 1$). Comparado con $\tau_g\approx\SI{37.5}{\uacc}$,
        cada unidad del accionamiento equivaldría a unos \SI{7}{\milli\newton\metre}. Es sólo una
        estimación de orden de magnitud; sirve para contrastar con el manual del accionamiento
        \verificar{escala del registro 16387 y par nominal del motor}.
\end{itemize}

\section{Limitaciones de las medidas}
\label{sec:res-limitaciones}

\begin{itemize}
  \item \textbf{Calibración de la célula.} La constante de \num{14.41} cuentas/N y el valor de
        \SI{51.78}{\newton} en la tara son provisionales: no proceden de una calibración con masas
        patrón. Las variaciones relativas del tramo central dependen poco de ellos, pero los valores
        absolutos en newtons deben tomarse como orientativos. Se propone el procedimiento de
        calibración del apartado siguiente.
  \item \textbf{Filtro de la célula.} La media móvil de ocho lecturas válidas introduce un retardo de
        varias décimas de segundo. Promediar ambos sentidos lo compensa en primera aproximación, pero
        limita el análisis de transitorios.
  \item \textbf{Lectura bruta por lote.} La lectura sin filtrar sólo se transmite una vez cada
        \SI{200}{\milli\second}.
  \item \textbf{Posición integrada.} No se ha contrastado con una medida independiente (codificador
        del accionamiento o medida angular de la barra).
  \item \textbf{Unidades del par.} El par se usa en unidades internas del accionamiento.
  \item \textbf{Una sola masa y un solo prototipo.} No se ha estudiado la dependencia con la carga
        ni con el montaje.
\end{itemize}

\subsection{Procedimiento de calibración propuesto}

\begin{enumerate}
  \item Colgar la célula en vertical con su cadena de montaje y tarar sin carga.
  \item Aplicar masas patrón en escalones crecientes y decrecientes que cubran el rango de uso (por
        ejemplo, de \num{0} a \SI{6}{\kilo\gram} en pasos de \SI{0.5}{\kilo\gram}), registrando
        \SI{10}{\second} por escalón con el eje parado.
  \item Ajustar por mínimos cuadrados la recta cuentas--fuerza y obtener la sensibilidad, el cero,
        el error de linealidad, la histéresis entre carga y descarga y la repetibilidad en tres
        series.
  \item Estimar la incertidumbre combinada, con las contribuciones de las masas patrón, la
        resolución, el ruido, la linealidad y la histéresis, y expresarla con $k=2$.
  \item Introducir la nueva constante en la interfaz y recalcular los ensayos existentes, cuyas
        cuentas brutas se conservan.
\end{enumerate}
